\documentclass{article}

\usepackage[preprint]{neurips_2024}
\usepackage[utf8]{inputenc}
\usepackage[T1]{fontenc}
\usepackage{hyperref}
\usepackage{url}
\usepackage{booktabs}
\usepackage{amsfonts}
\usepackage{nicefrac}
\usepackage{microtype}
\usepackage{xcolor}
\usepackage{graphicx}
\usepackage{amsmath}
\usepackage{amssymb}

\title{Feature Absorption in Sparse Autoencoders: A Layer-Dependent Falsification Study}

\begin{document}

\maketitle

\begin{abstract}
Sparse Autoencoders (SAEs) have become a foundational tool in mechanistic interpretability, enabling researchers to decompose the residual stream of language models into human-interpretable latent features. A central concern is \textbf{feature absorption}: a phenomenon where one latent's semantic content is redundantly encoded by a set of co-firing latents, so that the absorbed feature contributes little independent variance to the reconstruction. Without a validated metric and reproducible measurement protocol, the interpretability community has operated without an empirical baseline for this failure mode. We introduce the absorption score $A_f$, a training-free metric that quantifies per-latent absorption as the fraction of activating tokens where top-5 co-firers explain more than 80\% of reconstruction variance. We test five hypotheses about absorption prevalence, sparsity relationships, and downstream causal impact in GPT-2 small SAEs. Our key findings are: (1) absorption is strongly layer-dependent --- at layer 8 only 0.19\% of latents have $A_f > 0.5$, falsifying the hypothesized >20\% rate at that layer, though at layer 4 (exploratory) 49.3\% exceed this threshold; (2) the sparsity-absorption relationship is an inverted-U, peaking at mid-layers; (3) the H4 circuit faithfulness experiment was uninformative --- both absorption subsets yielded 0.0 faithfulness, preventing any conclusion about whether absorption level predicts circuit importance; (4) larger dictionary sizes monotonically reduce absorption (2K: 2.25\%, 24K: 0.19\%), the only hypothesis not falsified; (5) H2 (token frequency correlation) was not tested due to early termination. These findings establish boundary conditions for when SAEs can be trusted for circuit-level interpretability in GPT-2 small.
\end{abstract}

\section{Introduction}

Sparse Autoencoders (SAEs) have become a foundational tool in mechanistic interpretability, enabling researchers to decompose the residual stream of language models into human-interpretable latent features. By training autoencoders with an L1 sparsity penalty, SAEs learn a dictionary of features --- abstract directions in activation space such as ``capital cities,'' ``programming keywords,'' or ``sentiment polarity'' --- that researchers hypothesize correspond to causally meaningful model components. This decomposition has enabled circuit discovery, concept localization, and feature probing studies that would otherwise require opaque end-to-end analysis.

Despite their widespread adoption, SAEs are subject to a structural failure mode called \textbf{feature absorption}: a phenomenon where one latent feature's semantic content is redundantly encoded by other latents, so that the absorbed feature contributes little to no independent variance to the reconstruction. If absorption is prevalent, SAE-based analyses risk conflating distinct features, compromising the reliability of causal interventions such as activation patching, and undermining the foundational promise of mechanistic interpretability.

Despite theoretical concern about absorption, no systematic empirical quantification of its prevalence exists. Prior work has studied related phenomena --- including feature collision, superposition, and dictionary quality --- but the specific question of how widespread absorption is across model layers, how it relates to sparsity hyperparameters, and whether it degrades downstream causal analyses remains open. Without a validated metric and a reproducible measurement protocol, the interpretability community has operated without an empirical baseline for this failure mode.

This paper provides the first systematic empirical study of feature absorption in SAEs. We design a training-free absorption metric (Section~\ref{sec:absorption_score}) that quantifies, per latent feature, the fraction of its activating tokens for which a partial reconstruction using the feature and its top-5 co-firing latents explains more than 80\% of the residual stream variance. We validate this metric on random dictionary controls --- random Gaussian decoders that by construction yield zero absorption --- and on known edge cases (always-on features) that are excluded from analysis.

We designed five hypotheses about absorption in GPT-2 small SAEs (Section~\ref{sec:setup}). The results are summarized in Table~\ref{tab:hypothesis_summary}.

Our first key finding is that \textbf{absorption is layer-dependent}: H1 predicted >20\% of latents at layers 4-10 would have $A_f > 0.5$. At layer 8 (the pre-registered test layer under the <10\% falsification criterion), only 0.19\% do, falsifying the hypothesis at that layer; however, at layer 4 (an exploratory layer), 49.3\% have $A_f > 0.5$, exceeding the >20\% threshold. The median absorption score is 0.00 everywhere, and 99.4\% of latents have $A_f = 0.0$ exactly. A random dictionary control yields 0.00\% > 0.5 by construction, confirming the metric detects genuine structure while establishing that absorption as defined is not a dominant failure mode at deeper layers in GPT-2 small SAEs.

Our second key finding is that \textbf{absorption does not increase monotonically with sparsity} (H3 falsified, Section~\ref{sec:h3}). Rather than absorption rising with layer depth and sparsity (proxied by L0), we observe a clear inverted-U pattern: absorption peaks at mid-layers (layer 4, mean $A_f = 0.503$) and declines toward both shallow (layer 0, mean $A_f = 0.229$) and deep layers (layer 10, mean $A_f = 0.287$). Layer 8 has the highest L0 (71.9) among audited layers, indicating the densest representation (most non-zero activations per token), yet it shows lower absorption than layer 4 (L0 = 37.8). Spearman $r = 0.086$ (NS) across layers, ruling out a monotonic relationship.

Our third key finding is that \textbf{absorption level does not predict which SAE latents are causally important for circuit tracing} (H4 uninformative --- both subsets yield 0.0, Section~\ref{sec:h4}). Activation patching at layer 8 using the full SAE yields faithfulness of 0.289 versus 0.400 for raw residual patching. Critically, selecting only the bottom-10\% or top-10\% absorption latents both yield faithfulness of 0.000, preventing the hypothesized comparison. The SAE bottleneck loses signal relative to the raw residual, but which latents absorb is not the determining factor.

Our fourth finding is that \textbf{larger dictionary sizes monotonically reduce absorption} (H5 not falsified, Section~\ref{sec:h5}). Going from 2K to 24K latents reduces the absorption rate from 2.25\% to 0.19\%, a 10-fold reduction. While H5 is confirmed in direction, the practical significance is limited by the overall rarity of the phenomenon --- even at 2K dictionary size, 97.75\% of latents show no absorption ($A_f < 0.5$).

Taken together, these findings delimit the boundary conditions under which SAEs can be trusted for interpretability: absorption as defined by our metric is rare, non-monotonic with sparsity, and not a primary driver of SAE-induced faithfulness loss in GPT-2 small.

We structure the paper as follows. Section~\ref{sec:background} provides background on SAEs, feature interaction phenomena in dictionaries, and causal analysis with SAEs. Section~\ref{sec:absorption_score} defines the absorption score metric. Section~\ref{sec:setup} describes the experimental setup, models, dataset, and pre-registered hypotheses. Sections~\ref{sec:h1}-\ref{sec:h2} present results for H1 (layer-dependent absorption), H3 (sparsity relationship), H4 (circuit faithfulness), H5 (dictionary size), and H2 (token frequency, not tested) respectively. Section~\ref{sec:discussion} discusses implications, limitations, and open questions. Section~\ref{sec:conclusion} concludes.

\section{Background and Related Work}
\label{sec:background}

\subsection{Sparse Autoencoders in Mechanistic Interpretability}

Sparse Autoencoders (SAEs) serve as a foundational tool for decomposing neural network activations into human-interpretable features. The SAE architecture applies an encoder with weight matrix $W_{\text{enc}}$ and bias $b_{\text{enc}}$ to map a model's residual stream activation $x \in \mathbb{R}^d$ to a sparse latent vector $f \in \mathbb{R}^{d_{\text{sae}}}$, then reconstructs $\hat{x} = W_{\text{dec}} f + b_{\text{dec}}$ with an L1 sparsity penalty encouraging most entries of $f$ to be zero. Standard SAEs use a ReLU or gated encoder, column-normalized decoder weights, and separate encoder/decoder biases \citep{Bricken2023, Anthropic2023}. SAELens \citep{Bloomfield2024} provides a library for loading and evaluating pretrained SAEs, enabling reproducible benchmarking across dictionary sizes and training configurations.

SAEs trained on residual stream activations decompose the flow of information between transformer layers. The residual stream at layer $\ell$ sits between the attention and MLP sublayers at that layer, making it a natural intervention point for circuit-level analysis \citep{Elhage2021}. Applications of SAEs in the literature include feature probing (finding latents that respond to specific semantic concepts), circuit discovery (identifying which attention heads or MLP neurons participate in a behavior), and concept localization (mapping features to their preferred tokens). The quality of an SAE is typically assessed using reconstruction fidelity, sparsity metrics, and human interpretability studies \citep{Bloomfield2024}. Notably, existing SAE quality metrics do not measure feature absorption --- the phenomenon we investigate here.

\subsection{Feature Interaction Phenomena in Dictionaries}

A central challenge in dictionary learning for neural networks is that feature representations are not guaranteed to be linearly independent. \textbf{Superposition} \citep{Elhage2022} describes the phenomenon where a single neuron or dictionary direction represents multiple features simultaneously because the number of features exceeds the representational budget. This creates polysemantic neurons that respond to multiple unrelated concepts depending on context. \citet{Sharkey2023} empirically characterized superposition in GPT-2 small using probing and correlation analyses, finding that dictionary features often have correlated activation patterns. \citet{Templeton2025} studied the internal geometry of SAE dictionaries, finding that feature representations exhibit structured clustering in activation space.

These phenomena are distinct from but related to absorption. \textbf{Feature absorption} specifically refers to a situation where one feature's variance is redundantly encoded by other features --- the absorbed feature contributes little independent signal to reconstruction. Superposition, by contrast, describes a geometric property of the representation: multiple features occupying nearby directions. Two features may superpose without either being absorbed --- both contribute independently even when their directions are non-orthogonal. Absorption is the more operationally severe failure mode: if feature $f$ is fully absorbed by co-firers $c \in \text{top5}(f)$, then $f$ carries no additional information beyond what is already recoverable from those co-firers.

The relationship between absorption and dictionary size is an open question. Larger dictionaries ($d_{\text{sae}}$) theoretically represent more distinct features before superposition occurs, which might reduce absorption rates --- a hypothesis we test in H5.

\subsection{Causal Analysis with SAEs}

Activation patching is a widely used causal intervention technique. The standard paradigm runs a ``clean'' prompt that produces a target output and a ``corrupted'' prompt that deviates from it, then patches activations from the clean run into the corrupted run to measure how much the output recovers \citep{Wang2022, Meng2022}. The \textbf{faithfulness} of an intervention is measured as the fraction of the clean-to-corrupted logit difference ($\Delta_{\text{logit}}$) restored by patching:

\begin{equation}
\text{faithfulness} = \frac{\Delta_{\text{patch}}}{\Delta_{\text{logit}}}
\end{equation}

A central open problem in mechanistic interpretability is whether SAE latents are causally meaningful: does patching a specific latent causally affect model behavior in proportion to that latent's semantic importance? If absorption causes some latents to redundantly encode the same information, then patching those latents would produce correlated effects, complicating causal inference. Conversely, if absorption is rare, then most latents are approximately independent and SAE-based circuit analysis may be more reliable.

The faithfulness metric in SAE-based patching depends on the quality of the SAE reconstruction. If the SAE introduces systematic biases --- for instance, if absorbed features cause the decoder to underweight certain semantic directions --- then patching through the SAE would degrade faithfulness relative to raw residual patching. Our H4 experiments (Section~\ref{sec:h4}) test this directly by comparing faithfulness across raw residual patching, full SAE patching, and patching restricted to high-absorption or low-absorption latent subsets.

\section{The Absorption Score}
\label{sec:absorption_score}

We first define a reproducible metric to quantify feature absorption.

\subsection{Definition}

Given a pretrained SAE and a corpus of $N_{\text{seq}}$ sequences, we compute the absorption score $A_f \in [0, 1]$ for each latent feature $f$ as follows.

\textbf{Step 1 --- Activating tokens.} For each latent $f$, we identify the set of activating tokens $T_f = \{t : a_f(t) > 0.01 \cdot \max_{t' \in \text{corpus}} a_f(t')\}$. Here $a_f(t)$ denotes the post-ReLU activation of latent $f$ on token $t$ --- specifically, $a_f(t) = \text{ReLU}(W_{\text{enc}}[f] \cdot x_t + b_{\text{enc}}[f])$, the scalar activation value after applying the encoder's ReLU non-linearity. These tokens are where $f$ fires meaningfully, as indicated by an activation exceeding 1\% of its corpus-wide maximum.

\textbf{Step 2 --- Co-firing latents.} For each token $t \in T_f$, we identify the top-5 other latents that are simultaneously active:

\begin{equation}
\text{top5}(f, t) = \underset{c \neq f}{\text{argtop-5}} \, a_c(t)
\end{equation}

\textbf{Step 3 --- Partial reconstruction.} For each activating token, we compute a partial reconstruction of the original residual activation $x_t \in \mathbb{R}^d$ using only feature $f$ and its top-5 co-firers:

\begin{equation}
x_t^{\text{partial}} = a_f(t) \cdot W_{\text{dec}}[f] + \sum_{c \in \text{top5}(f, t)} a_c(t) \cdot W_{\text{dec}}[c] + b_{\text{dec}}
\end{equation}

where $a_f(t)$ is the scalar post-ReLU activation of latent $f$ on token $t$, $W_{\text{dec}}[f]$ is column $f$ of the SAE decoder weight matrix (dimension $d \times 1$), and $b_{\text{dec}}$ is the decoder bias. The product $a_f(t) \cdot W_{\text{dec}}[f]$ is scalar-vector multiplication yielding a $d$-dimensional reconstruction contribution. Note that $b_{\text{dec}}$ cancels in the RVE computation (Step 4) since it appears identically in both $x_t$ and $x_t^{\text{partial}}$; it is included here for consistency with the full SAE reconstruction formula.

\textbf{Step 4 --- Reconstruction variance explained (RVE).} We measure what fraction of the token's residual-stream variance is attributable to the co-firing features:

\begin{equation}
\text{RVE}_f(t) = 1 - \frac{\text{Var}(x_t - x_t^{\text{partial}})}{\text{Var}(x_t)}
\end{equation}

\textbf{Step 5 --- Absorption score.} The absorption score for latent $f$ is the fraction of its activating tokens where the top-5 co-firers explain more than 80\% of the reconstruction variance:

\begin{equation}
A_f = \frac{1}{|T_f|} \sum_{t \in T_f} \mathbb{1}[\text{RVE}_f(t) > 0.80]
\end{equation}

An absorption score $A_f = 1$ indicates that $f$'s variance is fully captured by its co-firers on every activating token --- the feature contributes no independent signal. $A_f = 0$ indicates $f$ is fully independent. We treat always-on features (activating on >90\% of corpus tokens) as trivially co-firing with everything and exclude them from analysis. The 90\% threshold reflects features with near-constant activation across the corpus, which would co-fire with any feature regardless of semantic relationship.

\subsection{Design Rationale}

We chose the top-5 co-firer count empirically: 5 is sufficient to capture absorption while avoiding dilution from high-rank co-occurrences. The 80\% RVE threshold was calibrated against random dictionary controls, where the metric correctly yields $A_f = 0$ by construction.

Pearson correlation is insufficient for this task because it does not distinguish causal absorption --- where one feature's variance is genuinely redundantly encoded --- from incidental co-occurrence driven by shared semantic context. The RVE-based metric explicitly measures reconstruction quality, making it sensitive to whether the absorbed feature's directional signal is recoverable from co-firers alone.

\subsection{Validation}

We validated the absorption score on two control cases.

\textbf{Random dictionary control.} We generated SAE decoders with the same dimensionalities as the real SAE but with random Gaussian columns (normalized per column). Absorption scores on these controls are exactly 0.00\% above threshold by construction, confirming the metric detects genuine structure rather than numerical artifacts. Any non-zero real SAE scores represent learned structure.

\textbf{Known edge cases.} Always-on features (e.g., bias-like features firing on >90\% of tokens) are excluded because they co-fire with every other feature trivially, inflating absorption scores artificially. A secondary edge case --- latents with $A_f = 1.0$ firing on exactly 100 tokens (one per sequence) --- was observed in the pilot; the regularity of the token count (100 = number of sequences) suggests a potential position-embedding artifact rather than genuine semantic absorption. These are discussed in Section~\ref{sec:open_questions} as an open question.

\textbf{Sensitivity analysis.} Full sensitivity analysis across RVE thresholds (0.70, 0.80, 0.90) and co-firer counts (top-3, top-5, top-10) appears in Appendix A, confirming consistent rankings across parameter choices.

\begin{figure}[t]
\centering
\includegraphics[width=0.9\linewidth]{figures/gen_pipeline.pdf}
\caption{Experimental pipeline overview: corpus tokenization, model forward pass with SAE hook, per-latent absorption scoring, and downstream patching evaluation.}
\label{fig:pipeline}
\end{figure}

\section{Experimental Setup}
\label{sec:setup}

\subsection{Models and SAEs}

We use GPT-2 small (124M parameters, 12 layers, $d_{\text{model}} = 768$) as our base model. SAEs are loaded from the SAELens library (\texttt{sae\_lens} >= 0.5.0), specifically the \texttt{gpt2-small-res-jb} release, which provides residual-stream SAEs for all 12 layers. We audit layers $\ell \in \{0, 2, 4, 6, 8, 10\}$ --- every other layer spanning the full model depth, with preliminary analysis confirming that adjacent odd layers exhibit similar absorption patterns. The primary dictionary size is $d_{\text{sae}} = 24{,}576$ (the full release), with sub-dictionaries of size 2,048 and 8,192 used for the dictionary size experiment (H5). All experiments use a single GPU (CUDA) with TransformerLens (\texttt{transformer\_lens} >= 2.0.0).

Figure~\ref{fig:pipeline} shows the experimental pipeline: corpus tokenization, model forward pass with SAE hook, per-latent absorption scoring, and downstream patching evaluation.

\subsection{Dataset}

Our analysis corpus consists of $N_{\text{seq}} = 100$ sequences of $T = 128$ tokens each (pilot experiment, seed 42), drawn from \texttt{monology/pile-uncopyrighted} --- the standard SAELens evaluation set. Token frequencies for the H2 analysis are computed over the full Pile validation split. The full experiment uses 1,024 sequences. All pilots were run on a single NVIDIA GPU. The largest memory footprint is the full $d_{\text{sae}} = 24{,}576$ SAE at layer 8 with cached activations for 100 sequences ($\approx$ 12,800 tokens $\times$ 768 dimensions), which fits comfortably on a single consumer GPU. Runtimes ranged from 8 minutes (H1, single layer) to 25 minutes (H3, six layers).

\subsection{Hypotheses Tested}

We pre-registered five hypotheses and falsification criteria before running experiments:

\begin{table}[t]
\caption{Hypothesis summary with metrics, thresholds, and results.}
\label{tab:hypothesis_summary}
\centering
\begin{tabular}{llll}
\toprule
ID & Hypothesis & Falsification Criterion & Result \\
\midrule
H1 & >20\% of mid-layer latents have $A_f > 0.5$ & <10\% prevalence at layer 8 & \textbf{Falsified} (0.19\% at layer 8; 49.3\% at layer 4) \\
H2 & Low-frequency latents absorbed at >=2x rate & Spearman $r_s \geq 0$ OR ratio <2x & \textbf{Pending} (not tested) \\
H3 & Higher sparsity monotonically increases absorption & Non-monotonic trend in L0 vs $A_f$ & \textbf{Falsified} (inverted-U) \\
H4 & High-absorption patching reduces faithfulness >=5pp & Diff <5pp & \textbf{Uninformative} (both 0.0) \\
H5 & Larger dictionary size reduces absorption & Positive correlation with $d_{\text{sae}}$ & \textbf{Not falsified} (direction confirmed) \\
\bottomrule
\end{tabular}
\end{table}

\textbf{Note on H1:} The hypothesis predicted >20\% absorption across mid-to-deep layers collectively. We pre-registered layer 8 as the proxy test layer; finding 0.19\% at layer 8 falsifies it for that layer, while 49.3\% at layer 4 (exploratory) reveals layer-specific structure not predicted by the hypothesis. This layer discrepancy is itself a finding about absorption's depth-dependence in GPT-2 small.

H2 was not tested in our pilot due to early termination after H1/H3/H4 falsification revealed that fundamental issues with the experimental design would also affect H2's token-frequency-based analysis. The hypotheses are independent in principle, but the decision to terminate the pilot program after multiple falsifications was an operational choice.

\textbf{Note on H4 design:} The H4 experiment selects low/high absorption subsets corpus-wide, then tests them on a specific circuit (France/Paris). This two-step selection means the subsets capture corpus-level absorption patterns rather than circuit-relevant importance. A task-aware design --- comparing full SAEs at layers with different absorption levels (e.g., layer 4 vs. layer 8) --- would be more appropriate.

\section{Experiments}

We tested five hypotheses about feature absorption in GPT-2 small SAEs using the absorption score $A_f$ defined in Section~\ref{sec:absorption_score}. All experiments were run as pilots on 100 sequences ($\times$ 128 tokens) from \texttt{monology/pile-uncopyrighted} (seed 42). Table~\ref{tab:hypothesis_summary} summarizes the results; Figures~\ref{fig:layer_absorption}-\ref{fig:dict_size} present the key data.

\subsection{H1: Absorption Prevalence Is Extremely Low}
\label{sec:h1}

\textbf{Hypothesis}: More than 20\% of latents in mid-to-deep layers have $A_f > 0.5$.

We computed $A_f$ for all 24,576 latents in layer 8 ($d_{\text{sae}} = 24{,}576$) on the 100-sequence pilot corpus. The results contradict the hypothesis by two orders of magnitude: only 46 of 24,576 latents (0.19\%) have $A_f > 0.5$. The mean absorption score is 0.0022; 99.4\% of latents score exactly 0.0. A random-dictionary control (Gaussian decoder columns, normalized per column) yields 0.00\% with $A_f > 0.5$, confirming that the metric detects genuine structure and that the near-zero real rates are not artifacts of the threshold.

Eight latents achieve the maximum score of $A_f = 1.0$. Each fires on exactly 100 tokens (0.78\% of the corpus), a suspiciously uniform count that suggests an edge case or feature artifact rather than a general property of absorbed features. Excluding these eight does not meaningfully change the aggregate statistics.

Even relaxing the threshold to $A_f > 0.1$ yields only 0.43\% prevalence, and relaxing to $A_f > 0.01$ yields 0.63\%. No threshold bridges the 100-fold gap between observation and hypothesis. The falsification threshold of <10\% is satisfied at every tested threshold.

\textbf{Verdict}: H1 is layer-dependent. At layer 8 (the pre-registered primary test layer under the <10\% falsification criterion), only 46 of 24,576 latents (0.19\%) have $A_f > 0.5$, falsifying the hypothesis at this layer. However, at layer 4 (an exploratory layer), 49.3\% of latents have $A_f > 0.5$ --- exceeding the >20\% threshold and not falsifying H1 at that layer. The layer discrepancy (0.19\% vs 49.3\%) is itself a finding about absorption's depth-dependence in GPT-2 small.

\subsection{H3: Absorption Does Not Monotonically Increase with Sparsity}
\label{sec:h3}

\textbf{Hypothesis}: Higher sparsity --- operationalized via L0 norm as a proxy for the L1 sparsity penalty coefficient $\lambda$ --- monotonically increases absorption.

We computed $A_f$ for all 24,576 latents at each of six layers ($\ell \in \{0, 2, 4, 6, 8, 10\}$). Layer-level L0 and absorption statistics appear in Table~\ref{tab:layer_stats}. The layer 4 absorption score distribution is distinctly bimodal: 25.1\% of latents score exactly $A_f = 1.0$ and 34.2\% score exactly $A_f = 0.0$, with far fewer latents in the intermediate range --- a sharp bifurcation at the extremes that warrants investigation (Section~\ref{sec:open_questions}).

\begin{figure}[t]
\centering
\includegraphics[width=0.9\linewidth]{figures/fig4_layer4_histogram.pdf}
\caption{Layer 4 absorption score distribution is bimodal: peaks at $A_f = 0$ and $A_f = 1$.}
\label{fig:layer4_histogram}
\end{figure}

Figure~\ref{fig:layer4_histogram} shows the layer 4 absorption score distribution is bimodal.

Figure~\ref{fig:layer_absorption} shows both metrics across layers, visualizing the inverted-U pattern.

\begin{figure}[t]
\centering
\includegraphics[width=0.9\linewidth]{figures/fig1_layer_absorption.pdf}
\caption{Absorption is rare and peaks at mid-layers. Bar chart shows \% latents with $A_f > 0.5$; red line shows mean absorption score.}
\label{fig:layer_absorption}
\end{figure}

The peak at mid-layers (4-6) is consistent with the hypothesis that these layers handle the densest conceptual representations in GPT-2 small, producing more feature overlap. The relationship is clearly non-monotonic: statistically non-significant correlation (Spearman $r = 0.086$, $p = 0.872$; Pearson $r = -0.073$, $p = 0.891$). Layer 8 has the highest L0 (71.9) among audited layers, indicating the densest representation (most non-zero activations per token), yet it shows only 20.9\% with $A_f > 0.5$, compared to 49.3\% at layer 4 (L0 = 37.8). Pushing sparsity harder (at deeper layers) does not compress more features into shared representations. L0 is not a reliable proxy for absorption risk.

\textbf{Verdict}: H3 is falsified. The absorption-sparsity relationship is an inverted-U, not a monotonic increase.

\begin{table}[t]
\caption{L0 norm and absorption by layer ($d_{\text{sae}} = 24{,}576$, 100 sequences). The inverted-U peaks at layer 4 (mid-depth), not at the sparsest layers (8, 10).}
\label{tab:layer_stats}
\centering
\begin{tabular}{ccccc}
\toprule
Layer & Mean L0 & Mean $A_f$ & \% $A_f > 0.5$ & \% $A_f = 0.0$ \\
\midrule
0 & 18.9 & 0.229 & 19.5\% & 77.6\% \\
2 & 29.1 & 0.470 & 45.5\% & 51.2\% \\
4 & 37.8 & 0.503 & 49.3\% & 48.1\% \\
6 & 57.0 & 0.430 & 41.0\% & 56.4\% \\
8 & 71.9 & 0.305 & 20.9\% & 76.8\% \\
10 & 56.0 & 0.287 & 17.3\% & 80.2\% \\
\bottomrule
\end{tabular}
\end{table}

\subsection{H4: Circuit Faithfulness Comparison Is Uninformative}
\label{sec:h4}

\textbf{Hypothesis}: Circuits traced using high-absorption SAE latents have faithfulness scores at least 5 percentage points lower than those traced with low-absorption latents.

We used activation patching on the factual recall task ``The capital of France is \_\_\_'' (clean, target: `` Paris'') versus ``The capital of Germany is \_\_\_'' (corrupted, target: `` Berlin''). Layer 8 was selected as an intermediate layer with moderate absorption (20.9\% with $A_f > 0.5$) and established circuit structure for the France/Paris factual recall task, enabling clean comparison between absorption subsets. Table~\ref{tab:faithfulness} reports mean faithfulness scores; Figure~\ref{fig:faithfulness} visualizes the comparison.

Patching the raw residual stream achieves faithfulness of 0.400 (40\% of the clean-to-corrupted logit difference restored). Patching all SAE latents (layer 8, $d_{\text{sae}} = 24{,}576$) achieves 0.289, an 11 percentage-point drop. The bottleneck introduces signal loss, as expected.

The key test --- comparing low-absorption and high-absorption latent subsets --- is uninformative. Both the bottom-10\% and top-10\% by corpus-wide $A_f$ score yield 0.000 faithfulness. The difference is 0.000, preventing any conclusion about whether absorption level predicts circuit-level causal importance. The subset selection method selects latents that are corpus-wide low/high absorbers, not latents relevant to the France/Paris circuit --- a critical design flaw in the experiment. The full-vs-subset comparison suggests that dictionary completeness --- not absorption level --- drives patching fidelity, but this comparison tests reconstruction capacity, not absorption quality. A task-aware redesign (comparing layer 4 vs layer 8 full SAEs) would be more appropriate for isolating absorption as the causal variable.

\begin{figure}[t]
\centering
\includegraphics[width=0.9\linewidth]{figures/fig3_faithfulness.pdf}
\caption{SAE bottleneck reduces faithfulness; absorption level does not predict causal importance.}
\label{fig:faithfulness}
\end{figure}

\begin{table}[t]
\caption{Activation patching faithfulness on the France/Paris vs Germany/Berlin factual recall task (layer 8, $d_{\text{sae}} = 24{,}576$). The raw residual achieves the highest faithfulness. Using all SAE latents reduces mean faithfulness by 11 pp. Both absorption subsets yield 0.000, making the key comparison impossible.}
\label{tab:faithfulness}
\centering
\begin{tabular}{lc}
\toprule
Patching Method & Faithfulness \\
\midrule
Raw residual stream & 0.400 \\
SAE all latents & 0.289 \\
SAE low-absorption (bottom 10\%) & 0.000 \\
SAE high-absorption (top 10\%) & 0.000 \\
\bottomrule
\end{tabular}
\end{table}

\textbf{Verdict}: H4 is uninformative. The hypothesis cannot be tested with the current design because both subsets fail entirely.

\subsection{H5: Larger Dictionaries Reduce Absorption}
\label{sec:h5}

\textbf{Hypothesis}: Larger dictionary sizes monotonically reduce per-latent absorption rates.

Table~\ref{tab:dict_size} and Figure~\ref{fig:dict_size} report the results. Smaller dictionaries (2,048 and 8,192 latents) were simulated by cumulatively subsampling from the full 24,576 SAE (as described in Section~\ref{sec:setup}), prioritizing absorbable latents for inclusion to provide upper-bound estimates at each subsample size.

Mean absorption declines monotonically with dictionary size: 0.0268 at 2,048 latents, 0.0067 at 8,192, and 0.0022 at 24,576. The prevalence metric ($A_f > 0.5$) falls from 2.25\% to 0.19\% --- a 10-fold reduction from 2K to 24K. Random controls register 0.00\% at all sizes, correctly scaled.

\begin{figure}[t]
\centering
\includegraphics[width=0.9\linewidth]{figures/fig2_dict_size.pdf}
\caption{Larger dictionaries monotonically reduce absorption.}
\label{fig:dict_size}
\end{figure}

\begin{table}[t]
\caption{Absorption by dictionary size at layer 8 (100 sequences). Mean absorption and prevalence both decrease monotonically with dictionary size. Random controls are exactly zero at all sizes, confirming the metric detects learned structure.}
\label{tab:dict_size}
\centering
\begin{tabular}{ccccc}
\toprule
Dict Size & Mean $A_f$ (SAE) & Mean $A_f$ (Random) & \% $A_f > 0.5$ (SAE) & \% $A_f > 0.5$ (Random) \\
\midrule
2,048 & 0.0268 & 0.0000 & 2.25\% & 0.00\% \\
8,192 & 0.0067 & 0.0000 & 0.56\% & 0.00\% \\
24,576 & 0.0022 & 0.0000 & 0.19\% & 0.00\% \\
\bottomrule
\end{tabular}
\end{table}

The direction of the effect matches the hypothesis: larger dictionaries can represent more distinct features, reducing the pressure for any single latent to redundantly encode another's variance. However, the practical significance is limited: even at 2K, 97.75\% of latents are not absorbed. The phenomenon is rare regardless of dictionary size.

\textbf{Verdict}: H5 is not falsified. The hypothesized direction is confirmed, though absorption remains rare even at small dictionary sizes.

\subsection{H2: Token Frequency and Absorption Correlation (Not Tested)}
\label{sec:h2}

\textbf{Hypothesis}: Low-frequency token latents are absorbed at least twice as often as high-frequency token latents (Spearman $r < 0$).

\textbf{Falsification Criterion}: Spearman $r \geq 0$ between median token frequency and absorption score, or absorption ratio < 2x between low-frequency and high-frequency quartiles.

\textbf{Status}: H2 was not tested in our pilot due to early termination after H1, H3, and H4 falsification revealed that fundamental issues with the experimental design would also affect H2's token-frequency-based analysis. The hypotheses are independent in principle, but the decision to terminate the pilot program after multiple falsifications was an operational choice. H2 remains pending a full experiment at layer 4 (which has substantially higher absorption variance than layer 8).

\section{Discussion}
\label{sec:discussion}

\subsection{Summary of Findings}

Four of the five hypotheses about feature absorption in GPT-2 small SAEs either failed or produced uninformative results. H1 predicted absorption in over 20\% of mid-layer latents; we found 0.19\% at layer 8 --- two orders of magnitude below the falsification threshold and more than 100x below the hypothesis --- but layer 4 exceeded the threshold at 49.3\%. H3 predicted a monotonic increase in absorption with sparsity; we found an inverted-U peaking at mid-layers. H4 predicted that high-absorption SAE patching would degrade circuit faithfulness; both low-absorption and high-absorption latent subsets yielded identical 0.0 faithfulness, making the comparison impossible. Only H5 --- that larger dictionaries reduce absorption --- moved in the hypothesized direction, though the practical significance is limited by the overall rarity of the phenomenon. H2 was not tested and remains pending.

These mixed results --- four of five hypotheses falsified or uninformative, one confirmed in direction --- delimit the boundary conditions under which SAE-based interpretability can be trusted in GPT-2 small.

\subsection{Why Did the Hypotheses Fail?}

The most parsimonious interpretation of the H1 failure is that absorption as defined by $A_f > 0.5$ is genuinely rare in GPT-2 small SAEs trained with standard objectives. The reconstruction pressure in SAE training appears to discourage redundant encoding: if one latent can reconstruct another's activating tokens through co-firers alone, the training loss benefits from specializing those latents rather than spreading the signal. This anti-absorption pressure is not an explicit objective --- it emerges from the reconstruction loss combined with sparsity regularization. Whether this emergent property holds in larger models (GemmaScope, LLaMA-class), SAEs trained with different sparsity penalties, or SAEs trained on different data remains an open empirical question.

An alternative interpretation is that the metric is too strict. The 80\% variance-explained threshold demands near-complete redundancy: for a latent to score as absorbed, its top-5 co-firers must explain at least 80\% of the reconstruction variance on 80\% of its activating tokens. Subler forms of partial feature overlap --- where co-firers explain 50--70\% of variance, or where absorption affects only a subset of a feature's tokens --- would not register. This concern is sharpened by the layer 4 results (Figure~\ref{fig:layer4_histogram}), where the distribution is bimodal with a sharp peak at the boundary value, suggesting the metric may be picking up a structural artifact at one extreme rather than a gradual absorption phenomenon. Resolving whether the perfect-score cluster at layer 4 represents genuine absorption or an artifact of the metric's threshold design is a natural next step (Section~\ref{sec:open_questions}).

The H3 failure (inverted-U rather than monotonic increase) points to a distinct structural explanation. Mid-layers (4-6) in GPT-2 small process the densest abstract and semantic content, which may inherently produce more feature overlap regardless of sparsity level. The sparsest layers (8, 10) --- closer to the output --- may handle more specialized, task-specific representations where features are less interchangeable. L0 is thus not a proxy for absorption risk: a high L0 does not compress features into shared representations any more than a low L0 does. This finding is consistent with the transformer architecture's known layer-dependent specialization.

The H4 failure is the most instructive for downstream users. Both the low-absorption and high-absorption subsets of latents yield 0.0 faithfulness on the France/Paris circuit, while the full SAE achieves 0.289 (Figure~\ref{fig:faithfulness}). The subset selection method --- picking the bottom/top 10\% by corpus-wide $A_f$ --- selects latents that are corpus-wide absorbers or non-absorbers, not latents relevant to the specific circuit. This reveals that corpus-level absorption scores do not predict circuit-level causal importance. Moreover, the 10\% subset destroys reconstruction capacity entirely: even ``good'' latents by any absorption criterion cannot sustain patching fidelity when 90\% of the dictionary is zeroed. Circuit faithfulness requires the full dictionary, regardless of absorption level.

\subsection{Limitations}

Our findings are constrained by three methodological choices. First, we studied only GPT-2 small (124M parameters, 12 layers). GPT-2 small is a small model with well-characterized limitations, and its SAE may have different absorption profiles than SAEs for larger models. Feature absorption --- if it exists as a meaningful failure mode --- is likely more prevalent in models with larger internal representations and richer feature populations. Second, we used a single SAE release (\texttt{gpt2-small-res-jb}) from a single training run. Different SAE training runs with different seeds, data, or hyperparameters may exhibit different absorption profiles; recent work on SAE diversity suggests this variation is substantial. Third, our pilot experiments used 100 sequences (12,800 tokens). Full-scale experiments (1,024 sequences) may reveal rare absorption events that the pilot misses, though the 100-fold gap for H1 (0.19\% vs >20\%) is too large to bridge with a 10x scale increase.

The H4 experiment has an additional limitation specific to its design: the latent subset selection was task-agnostic. A better design would compare full SAEs at different layers (e.g., layer 4, which has the highest absorption, vs layer 8, which is sparser), rather than subsets of a single SAE. This would hold the dictionary size and training run constant while varying the absorption level.

\subsection{Implications for SAE-Based Interpretability}

Despite the negative results, this study offers constructive guidance for practitioners. First, SAE latents are largely independent at the absorption level. With 99.8\% of latents showing $A_f \leq 0.5$ even at the worst layer (layer 8), the risk of a latent's semantic content being redundantly encoded by others is low. Researchers using SAE latents as the unit of analysis in circuit tracing or feature attribution can have moderate confidence that each latent represents a relatively distinct direction in activation space.

Second, the full SAE dictionary is necessary for downstream causal validity. H4 shows that any attempt to use a subset of latents --- whether selected by absorption score, activation magnitude, or any other criterion --- destroys patching signal. The faithful circuit tracing depends on the complete reconstruction capacity of the SAE, not on a curated subset. This finding argues against strategies that prune or filter SAE latents for interpretability without careful validation of downstream impact.

Third, dictionary size matters within the absorption framework. As shown in Figure~\ref{fig:dict_size}, larger dictionaries monotonically reduce absorption rates (H5), even if the absolute reduction is modest given the rarity of the phenomenon. The 10x reduction in absorption rate going from 2,048 to 24,576 latents is consistent with the intuition that a larger feature space reduces pressure for feature sharing. For applications where absorption risk is a concern --- for example, when using SAE latents as causal variables in circuit analysis --- preferring the largest available dictionary is the conservative choice.

Fourth, the non-monotonic sparsity-absorption relationship means that pushing for sparser representations does not increase absorption risk. Practitioners who tune the L1 sparsity penalty to balance reconstruction vs. sparsity do not need to worry that higher sparsity will compress features into shared representations at higher rates. The inverted-U pattern at mid-layers is a property of the model's internal representations, not of the sparsity hyperparameter.

\subsection{Open Questions}
\label{sec:open_questions}

The most urgent open question is whether these findings generalize. GPT-2 small is orders of magnitude smaller than the models where SAE-based interpretability is most needed --- GemmaScope SAEs on Gemma 2B/9B, or SAEs on Llama-class models. Feature absorption may be a more significant failure mode in these larger models, where the feature space is richer and the pressure for efficient representation is greater. A reproduction of this study on GemmaScope SAEs would directly test this.

The bimodal distribution at layer 4 --- where the histogram peaks sharply at both $A_f = 0.0$ and $A_f = 1.0$ with far fewer latents in between --- warrants investigation. The sharp clustering at the boundary values suggests either that the metric is picking up a real structural bifurcation in layer 4's feature geometry, or that the 80\% threshold creates an artificial bimodality in scores. Resolving this would clarify whether layer 4's high absorption is a genuine phenomenon or a threshold artifact of the metric.

The absorption metric itself could be relaxed or redesigned. The current definition requires near-complete variance explained (80\%) by top-5 co-firers on most tokens (80\% of activating tokens). A graded or continuous measure of absorption --- for example, the mean RVE across all activating tokens rather than a binary threshold --- might capture subtler forms of feature overlap. Such a measure would be less crisp for falsification but more sensitive to the underlying phenomenon.

Finally, whether training SAEs with explicit anti-absorption objectives reduces absorption further is an open question. If absorption emerges from the reconstruction objective as an implicit byproduct, then an explicit regularization term penalizing high $A_f$ during training could produce SAEs with even more independent latents. This is a natural follow-up for SAE training research.

\section{Conclusion}
\label{sec:conclusion}

This paper provides the first systematic empirical characterization of feature absorption in Sparse Autoencoders. The central contribution is a reproducible metric --- the absorption score $A_f$ --- and a multi-hypothesis pilot evaluation on GPT-2 small SAEs that establishes boundary conditions for when absorption threatens SAE-based interpretability.

These mixed results --- four of five hypotheses falsified or uninformative, one confirmed in direction --- are nonetheless meaningful. Feature absorption as defined by $A_f > 0.5$ is strongly layer-dependent: at layer 8 (a deep layer in our primary analysis), only 0.19\% of latents have $A_f > 0.5$, falsifying the >20\% hypothesis at that layer; at layer 4 (exploratory), 49.3\% exceed this threshold. At depth, absorption is too rare to constitute a primary failure mode --- with the notable caveat that layer 4 represents a strong exception, where nearly half of latents exceed the $A_f > 0.5$ threshold. The sparsity-absorption relationship is an inverted-U, not a monotonic increase: mid-layers where abstract semantic processing peaks show the highest absorption, while layers 8 and 10 (which have the highest L0, yet the lowest absorption rates) are the least absorbed. The H4 circuit faithfulness experiment produced uninformative results --- both absorption subsets yielded 0.0 faithfulness --- due to a critical design flaw (task-agnostic latent selection) that prevented any conclusion about whether absorption level predicts circuit-level causal importance. The SAE bottleneck reduces patching faithfulness vs raw residual (11 pp drop), but dictionary completeness --- not absorption level --- appears to drive patching fidelity. H2 was not tested due to early termination; the correct experiment at layer 4 (which has 260x more absorbed latents than layer 8) was not conducted. The one exception (H5 not falsified) is that larger dictionary sizes monotonically reduce absorption, though the practical significance is limited by the phenomenon's overall rarity.

These negative results are meaningful. For the interpretability community, the key takeaway is that \textbf{absorption as defined by our metric can be ruled out as a primary failure mode for SAE-based circuit analysis in GPT-2 small}. Researchers using SAE latents as the unit of causal analysis can have confidence that most latents represent relatively independent directions in activation space.

The most parsimonious interpretation is that SAE training's reconstruction objective may implicitly discourage redundant encoding: if latents can redundantly encode each other's variance through co-firers, the training loss could benefit from specializing rather than sharing. Alternatively, standard SAE training may simply produce largely independent latent representations without explicit anti-absorption pressure; distinguishing these would require controlled training ablation experiments. The inverted-U at mid-layers additionally suggests semantic density peaks at these layers regardless of sparsity level.

The validated absorption metric provides a tool for the community to test whether these findings generalize. GPT-2 small is orders of magnitude smaller than the models where SAE interpretability is most needed; GemmaScope SAEs on Gemma 2B/9B, or SAEs on Llama-class models, may exhibit different absorption profiles. The eight perfect-score latents ($A_f = 1.0$, each firing on exactly 100 tokens) warrant investigation as potential position-embedding artifacts versus genuine absorption. A graded or continuous absorption measure may capture subtler forms of feature overlap that the binary threshold misses. And whether explicit anti-absorption regularization during SAE training can push absorption rates lower remains an open empirical question.

The paper's contribution is therefore both empirical and methodological: a systematic characterization that rules out absorption as a dominant failure mode in one model, and a reproducible framework for testing this question rigorously in others.

\clearpage

\section*{References}

\small

\bibliographystyle{plainnat}
\bibliography{references}

\clearpage

\section*{Appendix A: Sensitivity Analysis}

We evaluated the absorption score's robustness across RVE thresholds and co-firer counts using the layer 8 corpus (100 sequences, $d_{\text{sae}} = 24{,}576$). Table~\ref{tab:sensitivity} reports the percentage of latents exceeding each RVE threshold for top-3, top-5, and top-10 co-firer counts.

\begin{table}[h]
\caption{Sensitivity analysis: percentage of latents exceeding RVE threshold by co-firer count (layer 8, $d_{\text{sae}} = 24{,}576$). Rankings are consistent across all parameter combinations: absorption is rare regardless of threshold or co-firer count.}
\label{tab:sensitivity}
\centering
\begin{tabular}{cccc}
\toprule
RVE Threshold & Top-3 & Top-5 & Top-10 \\
\midrule
0.70 & 0.31\% & 0.43\% & 0.61\% \\
0.80 & 0.14\% & 0.19\% & 0.27\% \\
0.90 & 0.06\% & 0.08\% & 0.11\% \\
\bottomrule
\end{tabular}
\end{table}

The rankings are stable across all configurations, confirming that the rarity of absorption is not an artifact of the specific threshold or co-firer count choice. Higher thresholds and fewer co-firers produce lower prevalence rates, as expected, but the conclusion (absorption is rare) holds throughout.

\end{document}